\section{Aproximación de la función de verosimilitud: de lo ideal a lo factible}

El cálculo directo de $\nabla_{x_t} \log p(y | x_t)$ no es factible, ya que $p(y | x_t)$ involucra un mapeo complejo desde los datos intermedios con ruido $x_t$ hasta la observación $y$. Necesitamos una serie de aproximaciones para hacer viable el cálculo.

\subsection{Dificultades clave: $x_t$ vs. $x_0$}

Podemos escribir directamente la expresión para $p(y | x_0)$ —cuando $\mathcal{A}$ es una función lineal, $y | x_0 \sim \mathcal{N}(\mathcal{A} x_0, \sigma^2 I)$, lo cual constituye una distribución gaussiana sencilla.

Pero lo que necesitamos es $p(y | x_t)$, no $p(y | x_0)$. El problema radica en: \textbf{¿cómo establecer la conexión entre $x_t$ y $x_0$?}

\begin{warningbox}{No se puede sustituir directamente $p(y|x_t)$ por $p(y|x_0)$}
La aproximación más tosca consiste en asumir que $p(y|x_t) \approx p(y|x_0)$, es decir, tratar los datos intermedios con ruido como si fueran datos limpios para calcular. Esta es una aproximación terrible: si $p(x_0|x_t)$ fuera realmente una distribución delta (es decir, si se pudiera recuperar $x_0$ de manera determinista a partir de $x_t$), no tendríamos necesidad del proceso iterativo de eliminación de ruido en absoluto.
\end{warningbox}

\subsection{Fórmula de Tweedie: conexión entre $x_t$ y $x_0$}

La \textbf{fórmula de Tweedie} proporciona un método para estimar la media de $x_0$ a partir de $x_t$:

$$
\hat{x}_0 = \mathbb{E}[x_0 | x_t] = \frac{1}{\sqrt{\bar{\alpha}_t}} \left( x_t + (1 - \bar{\alpha}_t) \nabla_{x_t} \log p(x_t) \right)
$$

Donde $\nabla_{x_t} \log p(x_t)$ es exactamente el score predicho por el modelo preentrenado (o equivalentemente, obtenido a través de la predicción del ruido $\epsilon_\theta(x_t, t)$).

\begin{importantbox}{Función central de la fórmula de Tweedie}
La fórmula de Tweedie nos indica que, dado $x_t$, podemos calcular el \textbf{valor esperado posterior} de $x_0$, denotado como $\hat{x}_0$ (también llamado estimador de eliminación de ruido). Aunque esto representa solo la media y no la distribución completa, sirve como un puente clave para aproximar la función de verosimilitud.
\end{importantbox}

\subsection{Aproximación gaussiana: simplificación de $p(x_0 | x_t)$}

Para hacer viable el cálculo, hacemos una suposición aproximativa (aunque tosca):

$$
p(x_0 | x_t) \approx \mathcal{N}(\hat{x}_0, r_t^2 I)
$$

Es decir, aproximar la distribución posterior de $x_0$ dado $x_t$ como una distribución gaussiana con media dada por la estimación de Tweedie $\hat{x}_0$ y varianza escalar $r_t^2$.

\begin{warningbox}{Limitaciones de la aproximación gaussiana}
Esta aproximación es bastante tosca. Si $p(x_0|x_t)$ fuera realmente una distribución gaussiana, no necesitaríamos el proceso iterativo de eliminación de ruido; bastaría con muestrear $x_0$ directamente desde $x_t$. La verdadera $p(x_0|x_t)$ es una distribución multimodal mucho más compleja que la gaussiana. Sin embargo, sin entrenamiento adicional, esta es una de las mejores aproximaciones posibles.
\end{warningbox}

\subsection{Resumen de la sección}

A través de la fórmula de Tweedie y la aproximación gaussiana, hemos establecido un puente desde $x_t$ hasta $x_0$, transformando la $p(y|x_t)$ que no puede calcularse directamente en una distribución gaussiana susceptible de cálculo analítico. A continuación, derivaremos el algoritmo de guía específico basado en esto.



